\documentclass[11pt,a4paper]{article}
\usepackage[utf8]{inputenc}
\usepackage[T1]{fontenc}
\usepackage{amsfonts}
\usepackage[french]{babel}
\frenchbsetup{StandardLists=true} % à inclure si on utilise \usepackage[francais]{babel}
\usepackage{enumitem}
\usepackage{amssymb}
\usepackage{amsmath}
\usepackage[left=2cm,right=2cm,top=2cm,bottom=2cm]{geometry}
\usepackage{multirow}
\usepackage[dvipsnames]{xcolor}
\usepackage{gensymb} %symbole degré

\usepackage[T1]{fontenc}
\usepackage{fouriernc}
\usepackage{graphicx}

\usepackage{setspace}
\singlespacing

\usepackage{pstricks}
\usepackage{xkeyval}
\usepackage{pst-labo}
\usepackage{chemist} %chimie

\usepackage{tcolorbox}
\usepackage{multicol}
\usepackage{caption}
\usepackage{fancyhdr}
\pagestyle{fancy}
\renewcommand\headrulewidth{1pt}
\fancyhead[L]{Lycée Denis Diderot }
\fancyhead[R]{Classe de 1 Bac Pro }
\setcounter{page}{16}\fancyfoot[C]{page \thepage}
\fancyfoot[R]{Année 2024-2025}
\fancyfoot[L]{Alexis RUIZ}
\usepackage{multirow,tabularx,booktabs}
\usepackage[tikz]{bclogo} 

\newcommand{\cadreblanc}[1]{%
\medskip\framebox[\linewidth][c]{%
\rule{0mm}{#1mm}}\par
\medskip}

\usepackage{awesomebox}
\setlength{\aweboxvskip}{0mm}
\usepackage{pifont}
\usepackage{chemmacros}
\chemsetup{modules=all}
\setlength{\columnseprule}{.25mm}

\renewcommand{\arraystretch}{2.5}
\newcommand*\acocher{\quitvmode{\fboxsep0pt \fboxrule0.8pt \fbox{\vrule height1.8ex depth.1ex width0pt \vrule height0pt depth0pt width1.9ex }}}

\newcolumntype{C}[1]{>{\centering\arraybackslash }m{#1}}

\newcommand{\code}{\awesomebox[black]{2pt}{

\includegraphics[scale=0.15]{../Chapitre 1 - Energie et puissance électrique/images/frame.png}
}{black}}

\newcommand{\moteur}{\awesomebox[black]{2pt}{

\includegraphics[scale=1]{images/moteur.jpg} 
}{black}}

\newcommand{\defconclusion}{\awesomebox[black]{2pt}{\rotatebox{90}{\Large{\textbf{Conclusion}}}}{black}}

\newcommand{\defbox}{\awesomebox[black]{2pt}{
\includegraphics[scale=0.5]{../Chapitre 1 - Energie et puissance électrique/images/appel exam.png}}{black}}

  \newenvironment{itemize*}%  % listes compactes
    {\begin{termite}%
      \setlength{\itemsep}{0pt}%
      \setlength{\parskip}{0pt}}%
    {\end{itemize}}

\frenchbsetup{StandardLists=true}



\begin{document}


\begin{center}
 \LARGE{Exercices - partie 3 } \\[3mm]
\end{center}



\normalsize

\hrule\
\\[0,5cm]

  \raggedright
  
\fbox{\textbf{Exercice 1 }} QCM : Parmi les trois propositions, choisissez la bonne réponse. 

\begin{center}

    \begin{tabular}{|p{6cm}|C{3cm}|C{3cm}|C{3cm}|}
         \hline
       \ding{202} Quelle formule défini la quantité de matière $n$ ...& 
       $\Box  ~~n=\dfrac{m}{M}$ & 
       $\Box  ~~n=\dfrac{M}{n}$     & 
        $\Box  ~~n=m \times M$  \\
         \hline
         
      \ding{203} M(C) = 12 g/mol. Quelle est la masse de $n = 2 mol$ de carbone ?   & 
      $\Box $ ~~$m = 12g$ & 
      $\Box $ ~~$m = 24g$  & 
      $\Box $ ~~ $m = 6g$  \\
        
         \hline
         
        \ding{204} La masse molaire moléculaire représente la masse &
        $\Box $ ~~ d'une mole de molécules  &
        $\Box $ ~~ d'une molécule  & 
        $\Box $ ~~ d'une mole d'atomes  \\
         
         \hline
         \ding{205} Comment se calcule la masse molaire de l'eau oxygénée \ch{H2O2}& 
         $\Box $ ~~ M(H) + 2 $\times$ M(O) &
         $\Box $ ~~ 2 $\times$ M(H) + 2 $\times$ M(O)  & 
         $\Box $ ~~ M(H) + M(O) \\
         
         \hline
         
         \ding{206} Comment s'appelle un liquide dans lequel on dissout un solide &
         $\Box $ ~~ un soluté &
         $\Box $ ~~ un solvant & 
         $\Box $ ~~ une solution \\
         
         \hline
         
         \ding{207} Il y a 0,5 mol d'ions \chemform{OH^-} dans 200 mL d'une solution. Quelle est la concentration en quantité de matière &
         $\Box $ ~~ $c = 0,025$ mol/L  &
         $\Box $ ~~ $c = 2,5$ mol/L  &
         $\Box $ ~~  $c = 0,25$ mol/L  \\
         \hline
         %\ding{208} Un mouvement rectiligne uniformément ralenti peut être représenté par ...  & $\Box $ ~~  & $\Box $ ~~  &  $\Box $\\
         %\hline
         \ding{208} Pour préparer avec précision 100 mL d'une solution de quelle verrerie a-t-on besoin & 
         $\Box$ ~~ \psset{unit=0.5cm}\pstTubeEssais[glassType=erlen] &
         $\Box$ ~~ \psset{unit=0.5cm}\pstTubeEssais[glassType=becher]&
         $\Box$ ~~ \psset{unit=0.5cm}\pstTubeEssais[glassType=fioleJauge] \\
         
         \hline
         \ding{209} Réaliser le titrage d'une espèce chimique en solution ; c'est déterminer   &
          $\Box$ ~~  sa concentration en quantité de matière $c$&
          $\Box$ ~~ sa masse $m$ &
          $\Box$ ~~ son volume $V$\\
         
         \hline 
         
         \ding{210} Lors d'un titrage colorimétrique, l'équivalence peut-être déterminée par  &   
         $\Box$ ~~ un volume de solution titrante égale au volume de solution à titrer &
         $\Box$ ~~ un changement de couleur de la solution à titrer dans le bécher&
         $\Box$ ~~ un changement de couleur de la solution à titrante dans la burette \\
         
         \hline

         
        
    \end{tabular}
\end{center}
\vfill
\newpage


\fbox{\textbf{Exercice 2 }} \\[3mm]

Vous devez préparer 200 mL d’une solution d'acide chlorhydrique de concentration 0,01 mol/L en ions \ch{H3O+} à partir d’une solution de concentration 0,1 mol/L.

\begin{enumerate}
    \item Replacez, dans le bon ordre, les vignettes ci-dessous.
    
    \centering
    \includegraphics[]{images/ex2.png}

    \raggedright
    \item Quel volume faut-il prélever dans la solution mère ?
    
    \cadreblanc{20}
\end{enumerate}


    


\fbox{\textbf{Exercice 3 }} 

\begin{enumerate}
    \item On réalise le titrage d'un volume \chemform{V_1 = 0,05~L} d'une solution d'acide chlorhydrique, de concentration \chemform{C_1} en ions \ch{H_3O^{+}} inconnue, par une solution d'hydroxyde de sodium de concentration \chemform{C_2=0,1~mol/L} en ions \chemform{HO^-}. \'A l'équivalence : \chemform{V_2=10,5~mL}. Calculez \chemform{C_1} en mol/L. 

\cadreblanc{15}


\item La concentration \chemform{C_1 = 0,009~mol/L} en ions \chemform{Cl^-} dans une eau minérale, déterminée par titrage, confirme-t-elle l'indication sur l'étiquette : "Chlorures : 322 mg/L"

\cadreblanc{22}

\end{enumerate}



\end{document}